\documentclass{article}
\usepackage{amsmath}

\title{Consolidation of Pair Q and P}
\author{Ada}
\date{}

\begin{document}
\maketitle

\section*{The pair in two sentences}
Q is a survey and framework for learning-augmented algorithms that requires a typed prediction object, a task-defined error, prediction-cost accounting, and a proof across each component interface. P is MARS-RA, which asks a large multimodal model for pairwise contribution comparisons, fits Bradley--Terry scores, and uses their softmax as a potential for reward shaping in an open multi-agent reinforcement-learning system.

\section*{The connexion pursued}
The peers tested whether Q turns P into a consistency--robustness result, with comparison error playing the role of prediction error. They concluded that this generic relabeling overstates the evidence: the useful connexion is Q's interface discipline, applied to P's chain from comparisons to rank aggregation, task credit, potential shaping, and finite-training performance (ledger entries 1--6, 13, and 28).

\section*{What was found}
First, P's estimation claim stops at a latent comparator score. If \(\widehat{\mathbf c}_t\) is the fitted score, \(\mathbf c_t^*\) the Bradley--Terry latent score, and \(\mathbf v_t^*\) task-grounded credit, then the relevant decomposition is
\[
 \inf_{a\in\mathbf R}
 \frac{\lVert\widehat{\mathbf c}_t-(\mathbf v_t^*+a\mathbf 1)\rVert_2}{\sqrt n}
 \leq
 \inf_{a\in\mathbf R}
 \frac{\lVert\mathbf c_t^*-(\mathbf v_t^*+a\mathbf 1)\rVert_2}{\sqrt n}
 +\frac{\lVert\widehat{\mathbf c}_t-\mathbf c_t^*\rVert_2}{\sqrt n}.
\]
The infimum accounts for Bradley--Terry translation invariance. Repeated queries can reduce the second, sampling term, but not the first, semantic-bias term. P's Shapley statement removes that bias only by assuming \(\mathbf c_t^*=\mathbf v_t^*\), so it does not prove recovery of task credit (entries 4--6 and 20).

Second, ordinary episodic potential-based reward shaping does not yield the desired comparison-quality theorem. For a consistently indexed potential,
\[
 \sum_{t=0}^{T-1}\gamma^t
 \bigl(\gamma\psi_{t+1}-\psi_t\bigr)
 =-\psi_0+\gamma^T\psi_T.
\]
With P's zero terminal potential and a fixed initial state, this is policy-independent. Thus comparison quality cannot change the exactly optimized policy through this term; P's empirical effects must instead concern finite-training dynamics (entries 7--9 and 20).

Third, P's displayed open-team construction is not consistently typed. It fits \(\widehat{\mathbf c}_t\in\mathbf R^{|I^t|}\) while the active set \(I^t\) changes, then subtracts consecutive vector potentials and combines a vector potential with a scalar environmental reward. If a fixed population-indexed potential \(\Phi_t\in\mathbf R^n\) is used and every coordinate transition is paid, then
\[
 \sum_{t=0}^{T-1}\gamma^t
 \bigl(\gamma\Phi_{t+1,i}-\Phi_{t,i}\bigr)
 =-\Phi_{0,i}+\gamma^T\Phi_{T,i},
\]
which restores the usual invariance under suitable endpoint conditions. If shaping is paid only while agent \(i\) is active on \([\tau_i,\sigma_i)\), the residual is
\[
 -\gamma^{\tau_i}\Phi_{\tau_i,i}
 +\gamma^{\sigma_i}\Phi_{\sigma_i,i},
\]
and it may be policy-dependent. The reward recipient, coordinate map, and entry/exit convention must therefore be stated before a robustness claim is meaningful (entries 22--23, 26, 29, 32, and 34).

Finally, the notes identify a conditional cost frontier. If a repaired estimator gives credit error at most \(b+AK^{-1/2}\), a learner-specific finite-training sensitivity bound supplies a constant \(L\), and each query costs \(\kappa\), then the controlled part of the bound would be
\[
 Lb+LAK^{-1/2}+\kappa K,
 \qquad
 K^*=\left(\frac{LA}{2\kappa}\right)^{2/3}.
\]
This follows by differentiating with respect to positive continuous \(K\). It is a proposed target, not a theorem from Q or P, because the required sensitivity result is absent (entries 11, 18, 27--29, and 32).

\section*{Objections and their status}
The first objection is that P's Proposition 1 is not well posed under connectivity alone. With two agents and all observed outcomes favoring the first, the comparison graph is connected but the unregularized Bradley--Terry likelihood improves as the score gap diverges, so a finite maximum-likelihood estimate need not exist. A constrained or regularized estimator, a gauge, dynamic-range control, and graph-dependent curvature are needed; P's main and appendix rate normalizations are also inconsistent (entries 7, 10, 15, and 25). This objection was accepted but not repaired in the pair materials.

The second objection is that exact PBRS makes an optimal-return consistency--robustness curve vacuous. It was resolved only conditionally: telescoping is valid after a common codomain and boundary payments are defined, whereas active-interval-only payments leave entry/exit terms (entries 23, 26, 29, 32, and 34). The third objection is that LMM calls need not be independent Bradley--Terry samples; systematic, correlated, cyclic, or misspecified comparisons do not average away. This remains unresolved and motivates separating semantic bias from sampling variation (entries 10, 12, and 15).

\section*{Outside sources}
No source outside Q and P was used. The ledger and peer notes derive their claims by auditing Q and P; they introduce no independent empirical data or external theorem beyond the PBRS and Bradley--Terry arguments worked out there.

\section*{What remains open}
The material is thin on positive theory. It does not provide a valid corrected Bradley--Terry rate, an open-team PBRS theorem with complete reward-allocation semantics, or a finite-training MAPPO regret, sample-complexity, or sensitivity bound. In particular, \(L\) and therefore the query frontier remain conjectural (entries 25, 28, and 32).

\section*{Smallest next step}
The smallest decisive step is to state one fixed-population, agent-indexed shaping rule that specifies scalar-to-agent allocation and entry/exit payments, then prove by explicit telescoping whether it preserves policy ordering. This settles the pair-specific interface question without requiring the harder MAPPO analysis; if invariance holds, the next theorem must target finite-training performance rather than optimal return (entries 28, 32, and 34).

\end{document}
